\section{The SOR relaxation correction and safeguarded over-relaxation}
\label{sec:omega}

An intermediate discussion incorrectly attempted to extend
\cref{eq:signed-l1-drop} to the full SOR convergence interval $(0,2)$.  The
mistake is the sign of $1-\omega$.  The correct statement is piecewise.

\begin{lemma}[Correct universal one-step $\ell_1$ bound]
\label{lem:omega-l1}
For the update \cref{eq:q-sor-update},
\begin{equation}
  \norm{q_{\rm sc}^+}_1
  \leq
  \norm{q_{\rm sc}}_1-\delta_\alpha(\omega)\abs{q_{{\rm sc},u}},
  \label{eq:omega-l1-general}
\end{equation}
where
\begin{equation}
  \delta_\alpha(\omega)
  :=1-\abs{1-\omega}-\chi\omega
  =
  \begin{cases}
    \dfrac{2\alpha\omega}{1+\alpha},&0<\omega\leq1,\\[2mm]
    \dfrac{2(1+\alpha-\omega)}{1+\alpha},&1\leq\omega<2.
  \end{cases}
  \label{eq:delta-omega}
\end{equation}
Thus a graph-independent signed-$q_{\rm sc}$ decrease is guaranteed only for
$0<\omega<1+\alpha$.
\end{lemma}

\begin{proof}
Remove the old $u$ contribution, add the retained value
$(1-\omega)q_{{\rm sc},u}$, and bound the neighbor contribution by
$\chi\omega\abs{q_{{\rm sc},u}}$:
\[
  \norm{q_{\rm sc}^+}_1
  \leq
  \norm{q_{\rm sc}}_1-\abs{q_{{\rm sc},u}}
  +\abs{1-\omega}\abs{q_{{\rm sc},u}}
  +\chi\omega\abs{q_{{\rm sc},u}}.
\]
This is \cref{eq:omega-l1-general}; simplifying separately on either side of
$\omega=1$ gives \cref{eq:delta-omega}.
\end{proof}

\begin{proposition}[Sharpness of the universal range]
\label{prop:omega-counterexample}
Take a vector $q_{\rm sc}$ supported only at the updated coordinate $u$.  Then
\begin{equation}
  \frac{\norm{q_{\rm sc}^+}_1}{\norm{q_{\rm sc}}_1}
  =\abs{1-\omega}+\chi\omega.
  \label{eq:omega-counterexample}
\end{equation}
For $\omega>1+\alpha$, this ratio exceeds one.  Hence no universal signed
weighted-$\ell_1$ contraction theorem can hold on the entire interval
$1<\omega<2$.
\end{proposition}

The optimal parameter \cref{eq:omega-star} lies in the universal
$\ell_1$-safe interval only when
\begin{equation}
  \alpha\geq(\sqrt2-1)^2=3-2\sqrt2\approx0.1716.
  \label{eq:omega-safe-alpha}
\end{equation}
For the small-$\alpha$ regime where acceleration is most interesting,
$\omega_\star$ is generally not weighted-$\ell_1$ monotone.

\subsection{Safeguarded optimal SOR}

Over-relaxation can still be used without sacrificing the proof.  Before
committing the candidate $\omega_\star$ update at $u$, compute its exact local
change in $M=\norm{q_{\rm sc}}_1$:
\begin{align}
  \Delta_u(\omega_\star)
  :=&\ \abs{q_{{\rm sc},u}}+\sum_{j\sim u}\abs{q_{{\rm sc},j}}
  -\abs{(1-\omega_\star)q_{{\rm sc},u}}
  \nonumber\\
  &-\sum_{j\sim u}
     \abs{q_{{\rm sc},j}
          +\chi\omega_\star q_{{\rm sc},u}/d_u}.
  \label{eq:exact-local-mass-change}
\end{align}
For $\sigma\in(0,1]$, accept the candidate only if
\begin{equation}
  \Delta_u(\omega_\star)
  \geq
  \sigma\frac{2\alpha}{1+\alpha}\abs{q_{{\rm sc},u}}.
  \label{eq:safeguard-test}
\end{equation}
Otherwise execute the $\omega=1$ update.  The fallback always passes with
$\sigma=1$ by \cref{lem:signed-l1}.  Hence every committed update has the
same amortization as \cref{thm:mass-tail}, up to $1/\sigma$ and a constant
factor for evaluating the candidate.  This gives a proof-safe practical tail
that uses $\omega_\star$ whenever local cancellation is favorable.

\begin{remark}[Plain optimal SOR still converges]
The failure of weighted-$\ell_1$ monotonicity does not imply divergence.
\Cref{lem:sor-decrease} proves objective descent for every
$0<\omega<2$, including $\omega_\star$.  What is lost is the
$1/\epsppr$ residual-mass amortization.  The unconditional objective argument
alone gives an $\epsppr^{-2}$ tail bound unless the handoff gap is chosen as in
\cref{thm:objective-tail}.
\end{remark}

\section{How many AESP stages are needed before the switch?}
\label{sec:burnin-stages}

The outer rate gives an explicit answer for the objective-gap switch.

\begin{proposition}[AESP stage count to the objective handoff]
\label{prop:J-bound}
Let $J$ be the first integer satisfying \cref{eq:objective-switch}.  It is
sufficient to choose
\begin{equation}
  J+1
  \geq
  \frac{10}{9}\sqrt{\frac{1-\alpha}{\alpha}}
  \log_+\!\left(
    \frac{200(1-\alpha)(1+\alpha)^2}
         {\alpha^{5/2}\epsppr}
  \right),
  \label{eq:J-bound}
\end{equation}
where $\log_+(z):=\max\{0,\log z\}$.  Therefore
\begin{equation}
  J
  =\bigO\!\left(
    \frac1{\sqrt\alpha}\log\frac1{\alpha\epsppr}
  \right).
  \label{eq:J-order}
\end{equation}
\end{proposition}

\begin{proof}
By \cref{eq:aesp-outer-bound}, it suffices that
\[
  \frac{200(1-\alpha^2)}{\alpha}
  (1-\varpi_{\mathrm A})^{J+1}
  \leq\frac{\alpha^{3/2}\epsppr}{1+\alpha}.
\]
Rearranging and using
$(1-\varpi_{\mathrm A})^{J+1}\leq e^{-\varpi_{\mathrm A}(J+1)}$ with
$\varpi_{\mathrm A}=0.9\sqrt{\alpha/(1-\alpha)}$ proves
\cref{eq:J-bound}.
\end{proof}

The stage count is accelerated, but stage count is not degree-weighted work.
A local theorem must also bound the volume processed inside each shifted
solve.
